\chapter{Arcane Magic}
\label{ch:arcane}

Arcane Magic is an approach to magic that acknowledges that there are magical rules that govern the Universe and that by studying these rules a Magician can manipulate reality to his will.

Often Arcane Magic is atheistic, regarding gods and spirits as merely intelligent forces of the Universe, that exist to be interacted with and dealt with on an equal footing. 

Practitioners of Arcane Magic develop in one of two ways. The majority are organised into schools of wizardry, which have their own books of spells and rules that they teach their apprentices. Alternatively, there is a long tradition of practitioners working in solitude, cut off from other Arcane Magic users and society at large, to focus purely on their magical activities. Occasionally they take on an apprentice, to teach their art, or simply as a helping hand around the magical laboratory.

\section{Arcane Magic Ranks}
Some schools may have ranks but that will depend on the specific campaign. Regardless, the skill increase bestows certain benefits.
\begin{description}
\item[Apprentices:] Students of Arcane Magic who will only know a couple of spells at an Arcane Magic Casting skill of 26\% to 50\%. Apprentices can spend two Improvement Points to create Arcane Crystals. An Arcane Crystal focuses up to one spell that the character knows per ten points of Arcane Magic Casting. These spells can be changed in the downtime between adventures, and have a star placed against the spell name on the character sheet to show they are focused. Spells that are focused by the Arcane Crystal increase the character’s Arcane Magic Casting (+20\% when used to all casting tests). The character is aware of where the Crystal is at all times and the enchantments placed upon it make it indestructible to non-magical means.

\item[Mages:] Graduates of the schools of wizardry or equivalent. They will know between five and ten spells, and will have an Arcane Magic Casting skill ranging from 51\% to 90\%. Mages can choose to spend two Improvements Points to create an Arcane Power Store item that allows storing up to POW additional Power Points. It can be any item, like a staff, medallion, ring, sword, etc. The character is aware of where this item is at all times and the enchantments placed upon the Arcane Power Store makes it indestructible to non-magical means.

\item[Archmages:] Acknowledged masters of Arcane Magic. They will have at least ten spells and an Arcane Magic Casting skill of 90\%+. Archmages can choose to spend two additional Improvement Points to improve their Arcane Power Store item to store a maximum of POWx2 Power Points. An Archmage with Arcane Magic mastery can spend two Improvement Points to Transcend.
\end{description}

\begin{figure}[h]
\begin{center}
  \includegraphics[scale=0.24]{img/ai-images/archmage.png}
\end{center}
\end{figure}

\subsection{Transcendence}
The Archmage has such high knowledge of the laws of Arcane Magic that they are no longer part of the Real World. Their body changes from a physical one, based upon CON and SIZ with hit points, to one based on pure magical energy, based upon POW and power points.

\subsubsection{Requirements}
\begin{rpg-list}
\item Arcane Magic Casting at 101\%.
\item POW at maximum (21 for humans).
\item Must know 10 different Arcane Magic spells.
\end{rpg-list}

\subsubsection{Benefits}
\begin{rpg-list}
\item They are immune to disease and poison, unless they specifically target POW or power points, and instead of taking damage to hit points, they take them to power points.
\item Spells which heal a physical body, such as Heal, Treat Wounds, and Divine Heal, do not work on the Archmage’s new magical form. Instead, an Archmage will transfer stored power points, such as those held in their Arcane Power Store, back into their magical body. This takes one combat round and is an action that requires full concentration to complete.
\item Defensive spells that protect against physical harm, such as Damage Resistance, still do so. As does physical armour, although an Archmage is more likely to use magical means to protect themselves.
\item The character no longer ages physically and is considered immortal.
\end{rpg-list}

For many player characters, Transcending is the end of their story in-game, with the player retiring the character from play. If the player has ideas of how the Transcended character can continue, they should discuss that with their Gamemaster to work out further adventures. Most non-player character Transcended Masters choose to withdraw from the Real World, using their new magical body to explore Other Worlds and the other avenues of magical power that open up to them. They often maintain a link with the Real World, where favoured students can contact them for sentimental reasons.



\section{Spot Rules}

\subsection{Learning Arcane Magic}
Arcane Magic is governed by the Arcane Magic Casting skill. The base chance for Arcane Magic Casting is INT. Spells are learnt separately, but the Arcane Magic Casting skill determines the success for casting all Arcane Magic spells. Characters learn Arcane Magic from other characters who know the practice. It costs eight Improvement Points to get access to the Arcane Magic discipline. The character gains the Mystic Vision spell for free.

\subsection{Learning Spells}
Before a spell can be cast using Arcane Magic, the following process must be followed:
The character must first learn the spell through research. In order to learn a particular Arcane Magic spell, the caster must possess the spell in written form or be taught by a teacher. In game terms this means having access to a teacher who knows the spell or a book or scroll where it is written down. The player then spends three Improvement Points and writes the spell down on their character sheet. Once the Arcane Magic spell has been learned, the character will be ready to try casting it.



\subsection{Casting Spells}
A character must be able to gesture with his hands, and be able to chant, in order to cast a spell. Whenever a spell is cast using Arcane Magic, there will always be a sight and sound that nearby creatures can detect, be it a flash of light, a crack of thunder, or a shimmering in the air. The exact effects are up to the Gamemaster and Player to decide, but will automatically be detected by any creatures within 10 metres per Magnitude of the spell. 

Casting a Arcane Magic spell requires a successful skill test using the Arcane Magic Casting skill. If successful, the spell takes effect. If the casting test fails, the spell does not take effect. 

\subsubsection{Power Points}
All Arcane Magic spells cost a base of one Power Point to cast. Arcane Magic spells can be modifying as the caster wishes (if he has the appropriate Power Points). If a Manipulation effect is applied to a spell, each effect costs one Power Point to apply (see below). 

\begin{figure}[h]
\begin{center}
  \includegraphics[scale=0.24]{img/ai-images/spell-fumble.png}
\end{center}
\end{figure}

The result of the Arcane Magic casting test depends on its success:
\begin{description}
	\item[Success:] A number of Power Points are deducted from the spellcaster’s total, equal to the Manipulation effects Power Points plus one of the spell. The spell then takes effect.
	\item[Failure:] The spell does not take effect and the character loses one Power Point for the base spell cost.
	\item[Critical:] Any attempt to resist or counter the spell suffers -20\% penalty. Moreover, only the base cost of one Power Point is lost (not for any Manipulations).
	\item[Fumble:] The spell fails and the Arcane Magic user loses 1D6 Power Points in addition to normal Power Point loss.
\end{description}



\subsubsection{Casting Time}
No other Combat Action may be taken while casting a spell, though the character may slowly walk up to half their Movement. 

A spell takes effect at the end of its casting, which starts at the beginning of the Combat Round and ends on the caster's Combat Order. Note that while spellcasting, a character will draw possible attacks from enemies they are adjacent to during a Combat Round. 

Distractions or attacks on the spellcaster as he casts will either automatically ruin the spell (if the spellcaster is blinded or disarmed, or suffers a Major Wound) or require Persistence tests for them to maintain concentration on the spell. 

\subsubsection{Spell Manipulations}
Arcane Magic spells have three basic effects which can be manipulated by the caster: Magnitude, Duration, and Range.

Each effect has a default value which the spell can be cast at, costing one Power Point. The default value for the spell effects are listed in table~\ref{tab:arcane-manipulations}.

\begin{table}
\begin{center}
\caption{Arcane Magic Manipulations}
\label{tab:arcane-manipulations}
\begin{rpg-table}[|c|c|c|Y|]
        \hline
	\textbf{Power Points}  & \textbf{Magnitude} & \textbf{Duration} & \textbf{Range}\\
        \hline
	1 (Default) & 1 & 5 minutes & 10m\\
	+1 & 2 & 15 minutes & 20m\\
	+2 & 3 & 30 minutes & 40m\\
	+3 & 4 & 1 hour & 80m\\
	+4 & 5 & 2 hours & 160m\\
	+5 & 6 & 4 hours & 320m\\
	+6 & 7 & 12 hours & 640m\\
	+7 & 8 & 1 day & 1km\\
	+8 & 9 & 2 day & 2km\\
	+9 & 10 & 5 day & 5km\\
	+10 & 11 & 1 week & 10km\\
	+11 & 12 & 2 weeks & 20km\\
	+12 & 13 & 1 month & 50km\\
	+13 & 14 & 2 months & 100km\\
	+14 & 15 & 3 months & 200km\\
	+15 & 16 & 6 months & 500km\\
	+16 & 17 & 1 year & 1000km\\
	+17 & 18 & 2 years & 2000km\\
	+18 & 19 & 5 years & 5000km\\
	+19 & 20 & 10 years & 10000km\\
	\hline
\end{rpg-table}
\end{center}
\end{table}

The tens value of the caster’s Arcane Magic Casting skill determines the max number of additional Power Points that can spend on each of the manipulation types. 

\begin{rpg-examplebox}
Omar the Magnificent with a Arcane Magic Casting skill of 80\% can spend an additional 8 Power Points on manipulating each of the spell’s effects, in Magnitude, Duration and Range. That’s a manipulation of up to 8 levels for each effect, not 8 levels in total across all three effects.
\end{rpg-examplebox}

The decision of which effects to manipulate and how many extra Power Points are to be spent is made before the spell is cast.

\begin{rpg-examplebox}
Lura casts Damage Boosting on Rurik’s sword, and wants it to be at a magnitude of 4 for an hour.

She has an Arcane Magic Casting skill of 60\%, which means she can spend an additional six Power Points on manipulating any spell’s effects. Looking at the Manipulation table, Lura can comfortably manage a Magnitude of 4, for three additional Power Points, and can manage a duration of an hour with another three points. 

Lura’s player rolls the dice and compares the result against Lura’s casting skill of 60\% to see whether she successful casts the spell.

In fact Lura, can spend a maximum of six points on a magnitude of range 640m, another six on a duration of 12 hours and another 6 on a magnitude of 7, which is a total of 19 Power Points (18 for the manipulations and 1 for the spell itself).
\end{rpg-examplebox}


\subsubsection{Dismissing Spells}
In a single Combat Round, a caster can dismiss any Permanent spell(s) he has cast, as a free action. Ceasing to cast a Concentration spell is immediate and not an action. 


\subsubsection{Extra Power Points}
As you can probably work out from the example above, it is possible for an Arcane Magic user to cast a spell which needs more Power Points in its manipulated form than an Arcane Magic user will normally have. Arcane Magic users get round this by carrying either Power Point stores (see Magic spell Create Magic Store) or other artifacts that can store Power Points (e.g. from other disciplines).


\subsection{Spell Traits}
All Arcane Magic spells share the same basic qualities:

\begin{rpg-list}
\item Base Magnitude of one 
\item Duration of 5 minutes, and
\item Range of 10 metres.
\end{rpg-list}

Other traits used by spells are detailed below. 
\begin{description}
	\item[Concentration:] The spell’s effects will remain in place as long as the character concentrates on it. Concentrating on a spell is functionally identical to casting the spell, requiring the spell caster to continue to gesture with both arms, chant and ignore distractions. This trait overrides the normal Arcane Magic spell default Duration. 
	\item[Instant:] The spell’s effects take place instantly. The spell itself then disappears. This trait overrides the normal Arcane Magic spell default Duration. 
	\item[Permanent:] The spell’s effects remain in place until they are dispelled or dismissed. This trait overrides the normal Arcane Magic spell default Duration.
	\item[Resist (Dodge/Persistence/Resilience):] The spell’s intended effects do not succeed automatically. The target may make a Dodge, Persistence or Resilience test (as specified by the spell) in order to avoid the effect of the spell entirely. Note that Resist (Dodge) spells require the target to be able to use Reactions in order to Dodge. In the case of Area spells, the Resist (Dodge) trait requires the target to dive in order to mitigate the spell’s effect. 
	\item[Touch:] Touch spells require the character to actually touch his target for the spell to take effect, using an Unarmed skill test to make contact. The caster must remain in physical contact with the target for the entire casting.
\end{description}

\section{Spells}

%\begin{samepage}
\begin{rpg-spell}
{Animate (Substance)}
{Concentration}

This spell allows the caster to animate the substance indicated, up to one SIZ for every point of Magnitude. The caster can cause it to move about and interact clumsily (Movement of 1m per three points of Magnitude). 

The caster’s chance to have the animated object perform any physical skill successfully is equal to his own chance to perform that action halved (before any modifiers). If the appropriate Form/Set spell is cast immediately after this spell, the caster can perform much finer manipulation of the object. In this case, the animated object will use the caster’s full skill scores for physical activities. 

This spell can only be used on inanimate matter. 
\end{rpg-spell}
%\end{samepage}


%\begin{samepage}
\begin{rpg-spell}
{Cast Back}
{}

This protective spell offers a chance of sending hostile spells (Arcane Magic and Magic) back to the attacking spell caster. 

Cast Back only affects spells that target the user specifically and have the Resist trait. Such spells may affect the protected character normally, but if it is resisted, the spell is launched back at the person who cast it, as long as its Magnitude is not greater than the Cast Back’s Magnitude. 
\end{rpg-spell}
%\end{samepage}



%\begin{samepage}
\begin{rpg-spell}
{Create Familiar}
{Permanent}

This spell allows the caster to make a personal sacrifice in order to create a Familiar. The Nature of the Familiar is dependent on the setting, but the spell typically only works on a non-sentient animal. The spell takes one hour to perform. The typical cost of ingredients for the spell is 100 Gold Pieces and it requires 3 Improvement Points. 

\begin{rpg-list}
\item An Arcane Magic user can have up to one Familiar for every ten points of Power they have.
\item A Familiar uses the animal as the template for stats.
\item A Familiar gains sentience when created equal to 2D6 INT and Magic equal to 2D6 POW, unless the animal had a higher score, in which case that POW is used. 
\item A Familiar has a permanent mental connection with its master; with a range equal to 1/2 the caster’s POW in kilometres. Through this link the Familiar can allow the caster to ‘see’ through its perceptual abilities. The Arcane Magic user can cast spells on the Familiar as if touching it and the Familiar can cast any spells it knows on its master.
\item A Familiar's characteristics, skills and magic can be increased by the player spending Improvement Points, as they would for the character.
\item When a Familiar is destroyed or killed the Arcane Magic user loses 1D6 Hit Points from the trauma. Half the Improvement Points spent on the familiar return to the character's pool of Improvement Points.
\end{rpg-list}

\end{rpg-spell}
%\end{samepage}


%\begin{samepage}
\begin{rpg-spell}
{Damage Boosting}
{Touch}

This spell can be cast upon any weapon up to five ENC. Each point of Magnitude adds one point to the weapon’s damage (the basic spell will increase a hatchet from 1D6 damage to 1D6+1 damage, for instance). 
\end{rpg-spell}
%\end{samepage}


%\begin{samepage}
\begin{rpg-spell}
{Damage Resistance}
{Touch}

This spell protects the body of the recipient. Any incoming attack dealing damage equal to or less than the Magnitude of the spell is ignored. Any incoming attack dealing more damage than the Magnitude of Damage Resistance is unaffected and will deal its full damage as normal. Note that the protected character may still suffer from Knockback if applicable. 
\end{rpg-spell}
%\end{samepage}


%\begin{samepage}
\begin{rpg-spell}
{Diminish (Characteristic)}
{Resist(Persistence/Resilience), Touch}

There are actually seven Diminish spells, one for each Characteristic. The spell will temporarily apply a penalty to the specified Characteristic equal to the Magnitude of the spell. The penalty applied by this spell may not reduce a Characteristic below one and a creature must have the Characteristic in question to be affected by this spell. 

Diminish (STR, DEX, CON and SIZ) are resisted with Resilience. Diminish (INT, POW and CHA) are resisted with Persistence. 

Applying a penalty to POW does not reduce the character’s Power Points.

Note that not all uses of this spell are malicious. Thieves and others often value the timely use of a Diminish (SIZ) spell, as it can greatly enhance their ability to enter restricted areas. 
\end{rpg-spell}
%\end{samepage}


%\begin{samepage}
\begin{rpg-spell}
{Dominate (Species)}
{Resist(Persistence)}

This spell allows the caster to gain control over a creature belonging to a specific species. If the target fails to resist the spell, it must obey the commands of the caster for the duration of the spell. 

The controlled creature shares a telepathic link with the caster by which it can receive its orders. If the caster and the creature dominated do not share a common language, the caster can order it about by forming a mental image of the actions they wish the dominated creature to perform.
\end{rpg-spell}
%\end{samepage}


%\begin{samepage}
\begin{rpg-spell}
{Energy Projection (Type)}
{Ranged, Instant, Resist(Dodge)}

Energy is either projected as a beam or a ball towards the target, which can avoid the attack by Dodging.

If the spell takes effect the target takes damage equal to double the Magnitude of the spell. Physical Armour does not protect against the damage, but magical protection does. Types of energy that can be projected by this spell are Cold (Dark), Lightning, Heat (Fire), Shards of Rock (Earth), Windblast (Air).
\end{rpg-spell}
%\end{samepage}


%\begin{samepage}
\begin{rpg-spell}
{Enhance (Characteristic)}
{Touch}

There are actually seven Enhance spells, one for each Characteristic. Essentially the reverse of the Diminish spell, Enhance allows the caster to temporarily apply a bonus to the specified Characteristic equal to the Magnitude of the spell. A creature must have the Characteristic in question to be affected by this spell. 

Applying a bonus to POW does not increase the character’s Power Points. 
\end{rpg-spell}
%\end{samepage}

\begin{figure}[h]
\begin{center}
  \includegraphics[scale=0.24]{img/ai-images/human-flying.png}
\end{center}
\end{figure}

%\begin{samepage}
\begin{rpg-spell}
{Fly}
{Concentration, Resist (Persistence)}

Using this spell allows the caster (or whomever or whatever he targets with the spell) to fly.  The caster may levitate a number of objects or characters (the caster counting as one of these characters if he so wishes). 

A levitated character may not be Overloaded and must have a SIZ Characteristic which is lower than the caster’s POW characteristic. 

Objects must have an ENC lower than the caster’s POW characteristic. 

Characters or objects moved by this spell have a base Movement Rate of 6m. All objects and characters moved by this spell move at the spellcaster’s behest, not their own. 

Each point of the spell’s Magnitude may either be used to increase the target’s Movement by +2m or to target an additional object or character – but not both. An Arcane Magic user casting this spell at Magnitude 4 may fly himself with a Movement of 12m, fly himself and a friend with a Movement of 8m each, or fly himself and three friends with a Movement of 6m each.
\end{rpg-spell}
%\end{samepage}


%\begin{samepage}
\begin{rpg-spell}
{Form/Set (Substance)}
{Instant}

There are an unlimited number of Form/Set spells in existence, one for every substance imaginable, from steel to smoke to water. 

Each point of Magnitude allows the caster to shape one ENC of solid substance or one cubic metre of an ethereal substance (like darkness). The caster must be Familiar with the shape he is forming. 

When the caster has finished the forming process, the substance retains its shape. Rigid substances like steel will hold the form they had at the end of the spell, while more mutable substances like water will immediately lose their shape. 

This spell can be used to mend damage done to an object. The caster must form the entire object and must succeed at an appropriate Craft test. If successful he will restore the item to its original condition. 

This spell can only be used on inanimate substances. 
\end{rpg-spell}
%\end{samepage}


%\begin{samepage}
\begin{rpg-spell}
{Glow}
{}

This spell causes a glowing point of light to appear on a solid substance. At its base, the spell creates an area of light one metre in radius, giving off the same illumination as a candle. Each additional point of Magnitude increases the radius of effect by one metre. At Magnitude 3, the brightness of the spell increases to that of a flaming brand at its centre. At Magnitude 5, it increases to that of a campfire and at Magnitude 10 to that of a bonfire. 

This spell can be cast on an opponent’s eyes. If cast on a living being the spell also gains the Resist (Dodge) trait. If the target fails to resist it, he will suffer a penalty to all attack, parry and Dodge tests, as well as any skills relying upon vision, equal to five times the spell’s Magnitude, until the spell ends or is dispelled. 
\end{rpg-spell}
%\end{samepage}


%\begin{samepage}
\begin{rpg-spell}
{Hamper}
{Resist (Resilience)}

Each point of Magnitude of Hamper subtracts 1m from the Movement rate of the target. Every two points of Magnitude also subtracts 1 from the recipient’s Combat Order.
\end{rpg-spell}
%\end{samepage}


%\begin{samepage}
\begin{rpg-spell}
{Haste}
{}

Each point of Magnitude of Haste adds 1m to the Movement rate of the recipient. Every two points of Magnitude also adds +1 to the recipient’s Combat Order.
\end{rpg-spell}
%\end{samepage}


%\begin{samepage}
\begin{rpg-spell}
{Holdfast}
{Touch}

This spell causes two adjacent ten centimetre by ten centimetre surfaces (roughly the size of a man’s palm) to commingle into one. The basic bond has a STR of 1. Each additional point of Magnitude will either increase the STR of the bond by +1 or double the area affected. 

This spell can affect organic and inorganic substances. If the caster is attempting to bond a living being with this spell, the spell gains the Resist (Resilience) trait.
\end{rpg-spell}
%\end{samepage}


%\begin{samepage}
%\begin{rpg-spell}
%{Make Potion}
%{Permanent}
%
%Using the arts of alchemy, Sorcerers can store Sorcery spells in a liquid known as a Potion.
%
%All potions have an attached cost of 1 Gold Ducat per magnitude of spell in ingredients. Multiple spells may be stored in one potion, with the understanding that they will all be cast when the potion is drunk at the magnitude that the creator originally put into the potion.
%\end{rpg-spell}
%\end{samepage}


%\begin{samepage}
\begin{rpg-spell}
{Make Scroll}
{Permanent}

These are written items which store Arcane Magic Spells.

All scrolls have an attached cost of 2 Gold Pieces per magnitude of spell in ingredients for special inks/parchments, etc.

The resulting scroll is a one-use item which, upon a successful Arcane Magic Casting test, casts the spell(s) with any manipulations, at the magnitude that was cast on the scroll.

Alternatively, upon a successful Arcane Magic Casting the reader of the scroll can learn the spell by spending the appropriate number of Improvement Points.

Either way, upon a successful use of the scroll, the spell fades from the scroll. If the casting roll merely fails the spell remains, but the reader can not attempt to use the scroll until their Arcane Magic Casting skill increases. If the casting roll is fumbled the spell fades from the scroll, without any benefit to the reader.
\end{rpg-spell}
%\end{samepage}


%\begin{samepage}
\begin{rpg-spell}
{Make Spell Matrix}
{Permanent}

This spell creates items that store Arcane Magic spells. 
The enchanter pays 1 Improvement Point per spell stored in the matrix.

All spell matrices have an attached cost of 20 Gold Pieces per spell in special materials.

The wielder can cast and manipulate the spell at the skill of the original enchanter, using their own Power Points as fuel.

Spell matrices are mundane items in their own right and if the item is broken, then the spell is dispelled. However at the time of enchantment the enchanter can spend another Improvement Point to magically harden the item, making it indestructible to non-magical means. %doubling its Hit Points and Armour Points.
\end{rpg-spell}
%\end{samepage}


%\begin{samepage}
\begin{rpg-spell}
{Mirage}
{Touch}

This is the Arcane Magic user’s version of Illusion. This spell creates an illusion based on all five senses. The illusion will seem real and solid unless the person looking at it succeeds in a Perception test. If the viewer succeeds in a Perception test and the Illusion could usually cause damage if believed in, it can no longer cause damage to that character. As soon as a viewer disbelieves the illusion it becomes insubstantial and ghost-like to him.

The Size and the strength of the illusion is governed by the Magnitude. A Magnitude 1 Illusion creates objects or creatures up to SIZ 10 and viewers have +60\% bonus to their Perception modifier test to disbelieve the illusion. For every extra Magnitude the caster can choose to either increase the size by 10 or apply a -20\% penalty to the Perception modifier test. Thus, with 5 Magnitude a caster could create a creature up to SIZ 30 with a Perception modifier of +20\%.
	
A positive Perception modifier means that the illusion is slightly fuzzy and unreal round the edges while a negative modifier means that it is indistinguishable from the real thing.
\end{rpg-spell}
%\end{samepage}


%\begin{samepage}
\begin{rpg-spell}
{Mystic Vistion}
{Concentration}

This spell allows the recipient to literally see magic. By augmenting the recipient’s natural vision, the spell allows them to see a creature’s Power Points, as well as enchanted items with their own Power Points or spells. The recipient must be able to actually see the creature or object for this spell to work. Mystic Vision also allows a recipient to see into the Spirit World.

On a normal success the recipient of the spell will only know roughly how many Power Points an object or creature has (1–10, 11–20, 21–30 and so forth). On a critical they will know exactly. On a fumble the Gamemaster should give the player a misleading total.

By looking at a spell effect, a recipient of Mystic Vision will automatically be aware of its magical origin (Arcane Magic, Divine Magic or Magic). By increasing the Magnitude of Mystic Vision, the caster can learn more about what he is seeing. Compare the Magnitude of Mystic Vision to the Magnitude of any spell that the target is either casting or under the influence of. As long as Mystic Vision’s Magnitude exceeds the other spell’s, the caster will be able to precisely determine the effects of the perceived spell, and a mental image of who cast the spell (if it is not obvious). 

By looking at an enchanted item, a recipient of Mystic Vision will automatically be aware of its gross magical effects (such as the types of enchantment currently on the item). Each point of Magnitude of Mystic Vision will also determine either the invested POW (and therefore the relevant strength) of a particular enchantment or a particular condition laid upon an enchantment, at the Gamemaster’s choice.
\end{rpg-spell}
%\end{samepage}


%\begin{samepage}
\begin{rpg-spell}
{Neutralise Magic}
{Instant}

This spell allows a caster to neutralise other spells. Neutralise Magic will eliminate a combined Magnitude of spells equal to its own Magnitude, starting with the most powerful affecting the target. If it fails to eliminate the most powerful spell then it will instead target the second-most powerful spell. As soon as Neutralise Magic can no longer dismiss a target’s spells, because all the remaining spell’s Magnitudes are too high, its effects immediately end. 

Neutralise Magic can be fired as a Reaction, but only when another spell is cast within Neutralise Magic’s Range that the character wishes to counter. A successful Neutralise Magic disrupts the other spell and nullifies it. As long as Neutralise Magic’s Magnitude equals or exceeds the target spell’s Magnitude, the target spell is countered. 
\end{rpg-spell}
%\end{samepage}

\begin{figure}[h]
\begin{center}
  \includegraphics[scale=0.24]{img/ai-images/other-world-portal.png}
\end{center}
\end{figure}

%\begin{samepage}
\begin{rpg-spell}
{Other World Portal (Other World)}
{Instant}

This spell creates a portal to a named Other World. The Magnitude of the spell is the number of creatures (SIZ 12-18) who can use the portal simultaneously. The portal exists as long as the spell is in effect. When the spell’s duration is reached, the portal closes instantly. 

If the spell casting is fumbled, catastrophic events occur. Here are some example events but the creative Gamemaster is encouraged to create more:
\begin{rpg-list}
\item A malignant creature from that Other World goes through the portal and attacks the caster, in an attempt to close the portal. 
\item The caster and all within 10m of them are sucked through the portal, which then promptly closes. Worse the caster is so befuddled that he cannot remember this spell for D20+D4 hours.
\item The Other World, to which the portal is connected, invades the home reality in a 1D10 km diametre from the portal. The home reality protects itself by throwing up a magical barrier that lets things into the beachhead but not out. 
\end{rpg-list}
\end{rpg-spell}
%\end{samepage}


%\begin{samepage}
\begin{rpg-spell}
{Palsy}
{Resist (Resilience)}

If the caster is able to overcome his target with this spell, he can turn the victim’s own nervous system against them. The spell will paralyse the target, provided the spell’s Magnitude is greater than the quarter of target’s current Hit Points. 
\end{rpg-spell}
%\end{samepage}


%\begin{samepage}
\begin{rpg-spell}
{Poison Antidote}
{}

This spell counteracts an antidote to any poison. For the duration of the spell, it reduces any poison’s potency by its magnitude X 5. This is for the chance to resist the poison initially. If that resistance test is failed, it still reduces the poison’s effective potency by its magnitude X 5.   If cast on a person who is already affected by a poison, they get another chance to resist the poison.
\end{rpg-spell}
%\end{samepage}


%\begin{samepage}
\begin{rpg-spell}
{Protective Sphere}
{}

When completed, the Protective Sphere will create a sphere-shaped area of protection with a radius in metres equal to the spell’s Magnitude. If this spell is cast on the ground (or other immovable place) it cannot be moved. If cast on a vehicle (such as the bed of a wagon) or a person, it will move with the target. After the sphere has been completed any one or all of the following spells can be added to provide the defensive capacities of the Sphere during the duration of the Sphere: Damage Resistance, Spell Resistance, Spirit Resistance. The Sphere on its own provides no protection.

The Protective Sphere’s perimetre contains the benefits of its combined Resistance spell(s). The Protective Sphere only inhibits spells or attacks entering the circle from the outside – attacks or spells originating within the circle are unaffected. Thus a Protective Sphere against spirits would block out outside spirits but have no effect on those already inside its perimeter. A Protective Sphere against damage or spells would block out incoming attacks/spells, but have no effect on those attacks made within the sphere (including attacks targeting those outside the sphere). 
\end{rpg-spell}
%\end{samepage}


%\begin{samepage}
%\begin{rpg-spell}
%{Remove Disease}
%{}
%
%This spell purges the recipient of any disease that is currently afflicting them. Make an opposed roll of the spell’s magnitude vs. the disease’s potency. If successful the character is restored to full health, any penalties to characteristics are removed and any hit points lost to the effects of the disease are regained.
%\end{rpg-spell}
%\end{samepage}


%\begin{samepage}
\begin{rpg-spell}
{Regenerate}
{Concentration Special, Instant, Touch}

This spell causes a severed or maimed limb to regrow or reattach. Regenerate cannot return a character from the embrace of death. 

The Magnitude of the spell must equal or exceed the maximum Hit Points lost as a result of the Major Wound taken. This spell will cause a limb severed by a Major Wound to regrow or, if the detached limb is still present, for the limb to reattach itself to its stump. 

Regenerate takes a number of minutes equal to the target’s SIZ to reattach a limb, during which time the caster must maintain concentration on the spell. The Hit Points lost due to the Major Wound are recovered at the end of this period.
\end{rpg-spell}
%\end{samepage}


%\begin{samepage}
\begin{rpg-spell}
{(Sense) Projection}
{Concentration}

Each (Sense) Projection spell is a separate spell. Each spell encompasses one of the five base senses but there are also variants for any unusual sensory mechanism appropriate to the game world (such as sonar). 

This spell allows the caster to project one of his senses anywhere within the spell’s Range. The spell forms an invisible and intangible sensor, some ten centimetres across, which receives the specified type of sensory input and transmits it to the caster. The sensor can move a number of metres per Combat Round equal to the spell’s Magnitude at the caster’s direction and allows them to use his Perception skill through the sensor. 

Spells can be cast through the sensor of some Projection spells. For instance, ranged spells require Sight Projection, while touch spells require Touch Projection (and likely Sight Projection too, simply so the caster can find his target efficiently). 

Characters using Mystic Vision can see the sensor and attack it if they wish, though it is only vulnerable to magic. Magical weapons and spells employed against the sensor will not destroy it but will instead transfer their damage directly to the caster.
\end{rpg-spell}
%\end{samepage}


%\begin{samepage}
\begin{rpg-spell}
{Sense (Substance)}
{Concentration}

Eminently useful for finding valuables from afar, this spell has a variation for every substance imaginable. Sense (Substance) will cause all sources of the substance within range of the spell to glow an appropriate colour visible only to the caster – diamonds will gleam like ice, amber will shine like a camp fire and so on. Each point of this spell’s Magnitude allows it to penetrate one metre of rock, wood or dirt. If the source is concealed behind such a material, the surface nearest the caster will glow for a moment. The spell cannot penetrate refined metal, though it can penetrate ore. 
\end{rpg-spell}
%\end{samepage}


%\begin{samepage}
\begin{rpg-spell}
{Shapechange (Species) to (Species)}
{Resist (Resilience), Touch}

Each Shapechange spell is a separate spell. Of all spells with multiple variations, the Shapechange spell has the most, comprising a new spell for almost every combination of creatures imaginable. The spell only works on living things – the dead or inanimate cannot be shapechanged. 

The Magnitude of the spell must be equal to or greater than the average SIZ of both specified species. Thus changing a mouse (SIZ 1) into a newt (also SIZ 1) is Magnitude 1. Changing a mouse into a lion (SIZ 19) is Magnitude 19. 

If the spell is successful, the target will be biologically changed, gaining the STR, DEX, CON and SIZ of its new form. Its INT, POW and CHA are unchanged, and the target retains its memories and abilities (though it may be unable to use some of those abilities in its new form). 
\end{rpg-spell}
%\end{samepage}


%\begin{samepage}
\begin{rpg-spell}
{Skin of Life}
{Touch}

This spell protects the recipient from suffocation by air deprivation, due to such factors as drowning or the Smother spell. Each point of Magnitude will cover three points of SIZ – thus a Magnitude 4 Skin of Life spell would sustain a SIZ 12 creature. 
\end{rpg-spell}
%\end{samepage}


%\begin{samepage}
\begin{rpg-spell}
{Smother}
{Concentration, Resist (Resilience Special)}

If successful, this spell neutralises the air surrounding the target, making each breath stale and worthless, depriving it of oxygen. The caster must concentrate each round, in order to keep the spell operating. For the duration of the spell, the target will be unable to breathe, essentially drowning on dry land. 

When the spell begins, the target’s Resilience test determines whether it is able to gasp in one last breath before Smother cuts off the surrounding oxygen supply. If the target succeeds, it may hold its breath as normal. If it fails, it will start drowning in the next Combat Round. 

This spell can also be used to extinguish fires, as the flames will be starved of oxygen. At Magnitude 1, it can extinguish a Flame, Magnitude 2 a Large Flame, Magnitude 4 a Small Fire, Magnitude 7 a Large Fire and Magnitude 10 it will put out an Inferno. Smother has no effect on magical fire or on fire-based creatures. 
\end{rpg-spell}
%\end{samepage}


%\begin{samepage}
\begin{rpg-spell}
{Spell Resistance}
{}

This spell matches its Magnitude against the Magnitude of any incoming spell. If Spell Resistance’s Magnitude is greater than the incoming spell’s, then the incoming spell has no effect. If the incoming spell’s Magnitude is equal to or greater than the Magnitude of Spell Resistance, then the spell affects the target normally. 

Unlike many protective spells, Spell Resistance remains in place for the entirety of its Duration – spells that successfully breach the spell do not dispel it. However, it does not discriminate between incoming spells – a comrade attempting to magically heal the recipient of Spell Resistance must overcome it in order to successfully use a healing spell. 
\end{rpg-spell}
%\end{samepage}


%\begin{samepage}
%\begin{rpg-spell}
%{Spiritual Projection}
%{Touch}
%
%Spiritual Projection causes the recipient’s soul to leave its corporeal body and manifest in the Spirit World. The recipient’s corporeal body slumps into a catatonic state for the remainder of the spell.  While Spiritual Projection is sometimes used for scouting purposes (as the recipient’s spirit can pass through nearly any obstacle) it is usually used to combat or confront Spirit World denizens.
%
%The recipient’s body remains vulnerable for the Duration of the Spiritual Projection. The soul will always know the direction its host body lies in, and the rough range to it in metres, but it cannot use the body’s senses. It is quite possible for a wandering soul to have its body destroyed and not realise it until it returns from its sojourn. It is also possible for a wandering spirit to possess the host body, leaving the wandering soul trapped as a spirit, at least until the spell ends and the soul dies. It is for these reasons that most Arcane Magic users ensure that their bodies are in a safe stronghold and protected by Damage, Spell and Spirit Resistance, before casting Spiritual Projection.
%
%As a traveller in the Spirit World, the recipient automatically gains the effects of Mystic Vision for the duration of his Spiritual Journey. However, he will suffer a –40\% penalty to all Perception tests to spot mundane items or events in the real world while in the Spirit World. A recipient may not travel further from its host body than the Range of the spell and it can move at double their normal rate.
%
%While on the Spiritual Journey the recipient otherwise obeys all the other rules that a Shaman does (see chapter \ref{ch:shamanism}) as regards skill use as well as initiating Spirit Combat. They may opt to use their Arcane Magic skill to make attacks, if it is higher than their Persistence. For the duration of the Spirit Combat the Spirit Damage of the Character is increased from 1D4 to 1D6.
%\end{rpg-spell}
%\end{samepage}


%\begin{samepage}
\begin{rpg-spell}
{Spirit Resistance}
{}

This spell matches its Magnitude against the POW of any spirit that comes into contact with the recipient. If the recipient of the spell’s combined POW + Spirit Resistance’s Magnitude is greater than the spirit’s POW, the spirit cannot touch the recipient. 

A spirit unable to touch a recipient will not be able to personally attack or harm them, including through ranged attacks. A spell cast by a spirit at the recipient is blocked unless its Magnitude exceeds the Spirit Resistances’s Magnitude. 
\end{rpg-spell}
%\end{samepage}


%\begin{samepage}
\begin{rpg-spell}
{Summon (Other World Creature)}
{Resist (Persistence)}

This spell allows the caster to summon one Other World creature, per casting, to the mundane world. The creature is not automatically under the caster’s control. If the summoned creature succeeds its Persistence test, it is free of the caster's command and if so inclined may be hostile to the caster. Otherwise it acts as if under the influence of a Dominate spell, for the Duration of the spell. The Duration of the summon spell also determines how long the creature is trapped in the mundane world.

Examples of of spell variations: Demons, Elementals, Spirits and Undead.
\end{rpg-spell}
%\end{samepage}


%\begin{samepage}
\begin{rpg-spell}
{Tap (Characteristic}
{Concentration, Resist (Persistence), Touch}

There are actually seven Tap spells, one for each Characteristic. These devastating spells allow the caster to permanently strip a target of Characteristic points, transforming them into Power Points for his own use. 

The caster must make contact with the target, either physically or through Touch Projection, in order to Tap it – therefore the spell cannot be used on incorporeal creatures, such as spirits. 

Tap will only work if its Magnitude is equal to, or greater than the target’s specified Characteristic. Thus a Magnitude 6 Tap Strength spell would only work on targets with a STR of 6 or lower. 

The number of points Tapped by the spell are equal to 1D6 per Combat Round the Spell is applied to the Victim.

Characteristic points lost to Tap are lost permanently, though the victim can raise them again through normal means of increasing a Characteristic. Characteristics may be Tapped to 0, which usually involves the death of the victim. The exception being Charisma.

For each Characteristic point the caster Taps, he will gain one Power Point. The caster is limited in the number of Power Points he can gain through Tap; the spell can only increase his Power Points to double his normal limit. An Arcane Magic user may simply Tap a target and dissipate any gained Power Points. 

If the caster gains more Power Points through Tap than his normal maximum, they will disappear at the rate of one Power Point per minute once the spell finishes. 
\end{rpg-spell}
%\end{samepage}


%\begin{samepage}
\begin{rpg-spell}
{Teleport}
{Instant, Resist (Dodge)}

Teleport allows an Arcane Magic user to instantaneously move themselves, or a target, to anywhere within the range of the spell, as long as the destination can be directly observed (Sense Projection spells allow the Caster to ‘see’ locations beyond physical line of sight), there is solid footing and no object bars their arrival. If these conditions are not met, the spell automatically fails. The caster is able to teleport objects up to 3 points of SIZ per point of Magnitude.
\end{rpg-spell}
%\end{samepage}


%\begin{samepage}
\begin{rpg-spell}
{Time Travel (Time Period)}
{Instant}

This spell transports the caster and a number of creatures (of SIZ 12-18) equal to the Magnitude of the spell to a named Time era via a Time Tunnel that opens up and instantly sucks them through to their destination. The Duration of the spell is the time that the caster and group jumps forward or backwards through time. 

Casters usually have some knowledge about the time period they are travelling to, and use an Anchor, a landmark such as a bronze statue, that exists in both the original and destination time period. If they are travelling blind without such an Anchor, the casting roll is at -20\% and the effects of a fumbled roll are even more catastrophic than the examples below suggest. 

If the spell casting is failed, the caster and group still travels, but they end up in the wrong location (1D10 Km away from the Anchor point) and time (1D10 time units away, the length of the time unit depends on Duration, e.g. if the duration was in days, the time unit is days).

If the spell casting is fumbled catastrophic events occur. Here are some example events; the creative Gamemaster is encouraged to create more:
\begin{rpg-list}
\item A Guardian creature from an Other World emerges from the portal and attacks the caster, in an attempt to close the portal. 
\item The caster, and all within 10m of them, is sucked through the portal which then promptly closes. The caster is so befuddled that he cannot remember the spell for D20+D4 hours. 
\item As above, but the caster and party arrive in a completely different Time Era or even an Alternative Reality.
\end{rpg-list}

Arcane Magic users with this spell can change time freely without having to worry about unintentional butterfly effect changes, or any alterations in their own existence or memory from changing their past. However, too regular use is likely to lead to the catastrophic effects of a fumble.
\end{rpg-spell}
%\end{samepage}


%\begin{samepage}
\begin{rpg-spell}
{Treat Wounds}
{Instant, Touch}

This spell must be cast upon a wounded character. It dramatically accelerates the natural healing rate of the target.  For every point of Magnitude of this spell, the caster can repair one hit point per Combat Round, for the Duration of the spell. Treat Wounds cannot reattach or regrow a severed limb and will not work on any Major Wound effects. 
\end{rpg-spell}
%\end{samepage}


%\begin{samepage}
%\begin{rpg-spell}
%{Venom}
%{Resist (Resilience), Touch}
%
%This spell infuses the target’s body with a magical poison. The Potency of the poison is equal to the spell’s Magnitude x 5, takes effect instantly, and does damage equal to the magnitude per combat round for the spell’s Duration. The target may resist the poison with a Resilience test, as normal.
%\end{rpg-spell}
%\end{samepage}


\section{Creating Magic Items}
Most Magic Items found in play in a game of Fantasy Crux will have been created by the characters or Non Player Character magicians.

Although the spells that their characters use to create Magic Items are detailed in their respective spell lists, its worth going through them briefly to remind yourself what spell does what and how it is used.

\begin{description}
\item [Make Scrolls] This spell allows for the creation of magical writings. Scrolls can be used to transfer knowledge (spells) or directly to cast the written spell.
\item [Make Spell Matrix] Arcane Magic users can create items and store spells using this spell. Spells stored can then be used and manipulated using the wielder's Power Points.
\end{description}

Arcane Magic users also find very useful the Magic spell Create Power Point Store so that they can have a greater pool for manipulations.

\subsection{Identify a Magic Item}
There is no catch all “Detect Magical Properties” or “Know Magic item” skill in Fantsy Crux. This is quite deliberate, keeping with the general policy that such items are not the equivalent of Magical shotguns.  Some options are:

Consult a Sage or other magical expert. This option will cost the characters lots of money. Take a baseline of one hundred silvers per point of spell magnitude or some perilous quest that the character must do in return. Such experts are rare, because most high ranking Magicians have little time for magical research for others, and would be more interested in their own schemes. 

Detect Magic spells. This merely tells you the item is magical.  A critical casting may tell the caster how powerful the item is.

Trial and error. The character tries to find out the item’s use by experiment. Allow creative and imaginative plans to reveal partially what the item does.

Researching the myths and legends around the item.  This is the most certain way of finding out what a magic item does. Of course such myths may be obscure themselves, requiring a dangerous adventure to a long hidden repository of knowledge to find.

